\nuevaReceta{Asado al Horno (Carne o Pescado)}{1h 20 min}{receta:asado_versatil}
    
    \begin{minipage}[t]{0.35\textwidth}
        \section*{Ingredientes}
        \begin{itemize}[leftmargin=*]
            \item \textbf{Base:} Patatas, cebolla, tomate y pimiento rojo.
            \item \textbf{Grasa:} Mantequilla.
            \item \textbf{Líquido:} Vino blanco.
            \item \textbf{Proteína:} Entraña, costillar, salmón o lubina.
            \item \textbf{Sazón:} Hierbas provenzales, ajo, perejil, sal y limón.
            \item \textbf{Extra:} Especias barbacoa (opcional).
        \end{itemize}
    \end{minipage}
    \hfill
    \begin{minipage}[t]{0.6\textwidth}
        \section*{Preparación}
        \begin{enumerate}
            \item \textbf{Base:} Precalentar el horno a 200°C. Untar el recipiente con mantequilla para crear una capa antiadherente.
            \item \textbf{Cama de verduras:} Colocar rodajas de patata, cebolla y tomate. Añadir el pimiento en tiras, las especias, el ajo y el perejil.
            \item \textbf{Líquido:} Verter un poco de vino blanco (sin llegar a cubrir las patatas para que se asen y no se cuezan).
            \item \textbf{Si es CARNE:} Aderezar la pieza, colocar sobre las verduras y hornear todo junto. Vigilar y añadir agua/vino si se seca.
            \item \textbf{Si es PESCADO:} Hornear primero las verduras (50-60 min). Sacar, poner el pescado aderezado con limón y hornear solo 10-15 min más.
        \end{enumerate}
    \end{minipage}

    \vspace{1em}
    \begin{tcolorbox}[colback=orange!5, colframe=orange!50, title=Nota de tiempos, size=small]
        El pescado requiere mucho menos tiempo; si lo metes desde el principio con las patatas, quedará muy seco. ¡Respeta los tiempos!
    \end{tcolorbox}
\end{receta}